\section{\sys{} Design}
\label{sec:design}

\sys{} consists of three components:
\begin{itemize}
  \item \emph{Affine folding} moves shared scales and offsets from individual KV values to matrix-product inputs and results so that Tensor Cores consume operands built directly from packed quantized values, reducing the affine demand in \Cref{fig:ops-per-byte} (\Cref{sec:design-fold}).
  \item \emph{Bit-level Tensor Core operand construction} builds Tensor Core operands directly from packed quantized values, reducing the INT$\to$FP conversion demand in \Cref{fig:ops-per-byte} (\Cref{sec:design-unpack}).
  \item \emph{Warp-specialized GPU pipelines} overlap data movement, operand preparation, matrix products, and affine corrections on Hopper and Blackwell (\Cref{sec:design-pipelines}). 
\end{itemize}

\subsection{The affine folding mechanism}
\label{sec:design-fold}

\tikzset{
  mgrid/.style={draw=black!55, line width=0.5pt},
  mframe/.style={draw=black!70, line width=0.6pt},
  mname/.style={font=\footnotesize\sffamily, text=black!75},
  fcell/.style={font=\scriptsize, text=black!85, inner sep=0pt},
  axname/.style={font=\footnotesize\sffamily, text=black!65},
  axarrow/.style={-{Stealth[length=3.4pt, width=2.8pt]}, draw=black!55, line width=0.5pt},
  szarrow/.style={{Stealth[length=3.2pt, width=2.6pt]}-{Stealth[length=3.2pt, width=2.6pt]}, draw=ptred!85!black, line width=0.7pt},
  szlab/.style={font=\footnotesize\sffamily, text=ptred!70!black, fill=white, inner sep=0.6pt},
  caselab/.style={font=\footnotesize\sffamily, text=black!85, align=center},
  mode/.style={font=\footnotesize\sffamily\bfseries, text=ptblue!80!black, align=center},
  qbox/.style={draw=black!70, line width=0.5pt, rounded corners=2pt, minimum height=0.54cm, minimum width=0.72cm, inner xsep=3pt, align=center, font=\footnotesize\sffamily},
  quant/.style={qbox, double, double distance=1pt, fill=ptcyan!20},
  prod/.style={qbox, fill=white},
  res/.style={qbox, fill=ptgreen!16},
  flow/.style={-{Stealth[length=4.2pt, width=3.2pt]}, draw=black!75, line width=0.7pt},
  foldline/.style={-{Stealth[length=4pt, width=3pt]}, draw=ptblue, line width=0.8pt, dashed},
  grpout/.style={draw=ptred!85!black, line width=1.1pt, rounded corners=1.2pt},
  foldbig/.style={-{Stealth[length=5pt, width=4pt]}, draw=ptblue, line width=1.2pt},
  foldlab/.style={font=\footnotesize\sffamily, text=ptblue!80!black, fill=white, inner sep=1pt, align=center},
}

\providecommand{\foldmgridd}[5]{%
  \foreach \gi in {1,...,#4}{\ifnum\gi<#4
    \draw[mgrid] (#1+\gi*#5,#2) -- (#1+\gi*#5,#2+#3*#5);\fi}
  \foreach \gj in {1,...,#3}{\ifnum\gj<#3
    \draw[mgrid] (#1,#2+\gj*#5) -- (#1+#4*#5,#2+\gj*#5);\fi}
  \draw[mframe] (#1,#2) rectangle ++(#4*#5,#3*#5);}

\def\fca{0.55}
\def\fcb{0.34}
\def\fcp{0.24}
\def\fcc{0.57}
\def\fcd{1.22}
\def\frh{0.38}
\def\dxs{3.56}
\tikzset{
  fdlab/.style={font=\scriptsize\sffamily, text=black!75, align=center, inner sep=1pt},
}
\providecommand{\fgrid}[6]{%
  \foreach \gi in {1,...,#4}{\ifnum\gi<#4
    \draw[mgrid] ({#1+\gi*#5},#2) -- ({#1+\gi*#5},{#2+#3*#6});\fi}%
  \foreach \gj in {1,...,#3}{\ifnum\gj<#3
    \draw[mgrid] (#1,{#2+\gj*#6}) -- ({#1+#4*#5},{#2+\gj*#6});\fi}%
  \draw[mframe] (#1,#2) rectangle ({#1+#4*#5},{#2+#3*#6});}
\providecommand{\ftint}[4]{%
  \fill[ptblue!26] (#1,#2) rectangle ++(#3,#4);}
\providecommand{\fcolvals}[4]{%
  \foreach \t [count=\i from 0] in {#4}{%
    \node[fcell] at ({#1+0.5*#3},{#2-(\i+0.5)*\frh}) {\t};}}
\providecommand{\frowvals}[4]{%
  \foreach \t [count=\i from 0] in {#4}{%
    \node[fcell] at ({#1+(\i+0.5)*#3},{#2+0.5*\frh}) {\t};}}
\providecommand{\foldbody}{%
  \fgrid{-1.91}{0}{3}{3}{\fca}{\frh}%
  \fgrid{0}{1.74}{3}{3}{\fcb}{\frh}%
  \fgrid{0}{0}{3}{3}{\fcb}{\frh}%
  \node[mname, anchor=south west] at (-1.91,1.22) {$A$};%
  \node[mname, anchor=west] at (1.13,2.24) {$q$};%
  \node[mname, anchor=south] at (0.51,1.22) {$C$};}
\providecommand{\foldaxes}{%
  \draw[axarrow] (-0.13,2.86) -- (-0.13,1.76);%
  \node[axname, anchor=east] at (-0.21,2.31) {$K$};}
\providecommand{\foldafull}{%
  \frowvals{-1.91}{0.76}{\fca}{{$A_{11}$},{$A_{12}$},{$A_{13}$}}%
  \frowvals{-1.91}{0.38}{\fca}{{$A_{21}$},{$A_{22}$},{$A_{23}$}}%
  \frowvals{-1.91}{0.00}{\fca}{{$A_{31}$},{$A_{32}$},{$A_{33}$}}}
\begin{figure}[t]
\centering
\subfloat[Operand-fold scales a column of $A$ and adds the offset correction to every output column (\Cref{eq:operand-fold}).\label{fig:fold-operand}]{%
\begin{tikzpicture}[x=1cm, y=1cm, mname/.append style={inner sep=1pt},
                    caselab/.append style={inner sep=1pt}]
  \useasboundingbox (-1.94,3.30) rectangle (6.26,-0.64);
  \ftint{0}{2.50}{1.02}{\frh}
  \ftint{1.14}{2.50}{0.48}{\frh}
  \fgrid{1.14}{2.50}{1}{2}{\fcp}{\frh}
  \frowvals{1.14}{2.50}{\fcp}{{$s$},{$z$}}
  \foldbody
  \frowvals{0}{2.50}{\fcb}{{$q_1$},{$q_2$},{$q_3$}}
  \fcolvals{-1.91}{1.14}{\fca}{{$A_1$},{$A_2$},{$A_3$}}
  \foldaxes
  \draw[grpout] (-0.03,2.47) rectangle (1.65,2.91);
  \node[caselab, anchor=south] at (-0.13,3.00) {before fold};
  \begin{scope}[shift={(\dxs,0)}]
    \ftint{-1.91}{0}{\fca}{1.14}
    \ftint{1.42}{0}{\fcc}{1.14}
    \fgrid{1.42}{0}{3}{1}{\fcc}{\frh}
    \fcolvals{1.42}{1.14}{\fcc}{{$zA_1$},{$zA_2$},{$zA_3$}}
    \foldbody
    \frowvals{0}{2.50}{\fcb}{{$q_1$},{$q_2$},{$q_3$}}
    \fcolvals{-1.91}{1.14}{\fca}{{$sA_1$},{$sA_2$},{$sA_3$}}
    \draw[flow] (1.35,0.57) -- (1.10,0.57);
    \node[axname, anchor=south, inner sep=0.5pt] at (1.225,0.66) {$+$};
    \draw[grpout] (-1.94,-0.03) rectangle (-1.33,1.17);
    \draw[grpout] (1.39,-0.03) rectangle (2.02,1.17);
    \node[fdlab, anchor=north] at (-1.085,-0.08) {fold into\\operand $A$};
    \node[fdlab, anchor=north] at (1.705,-0.08) {output\\correction};
    \node[caselab, anchor=south] at (0.055,3.00) {after fold};
  \end{scope}
  \draw[foldbig] (1.14,0.57) -- (1.51,0.57);
\end{tikzpicture}}
\\[-2pt]
\subfloat[Acc-fold scales a partial output column and adds the offset-weighted row sums of $A$ (\Cref{eq:acc-fold}).\label{fig:fold-acc}]{%
\begin{tikzpicture}[x=1cm, y=1cm, mname/.append style={inner sep=1pt},
                    caselab/.append style={inner sep=1pt}]
  \useasboundingbox (-1.94,4.16) rectangle (6.26,-1.14);
  \ftint{0.68}{1.74}{\fcb}{1.14}
  \ftint{0.68}{3.00}{\fcb}{0.76}
  \fgrid{0.68}{3.00}{2}{1}{\fcb}{\frh}
  \fcolvals{0.68}{3.76}{\fcb}{{$s$},{$z$}}
  \foldbody
  \fcolvals{0.68}{2.88}{\fcb}{{$q_1$},{$q_2$},{$q_3$}}
  \foldafull
  \foldaxes
  \draw[grpout] (0.65,1.71) rectangle (1.05,3.79);
  \node[caselab, anchor=south] at (-0.345,3.88) {before fold};
  \begin{scope}[shift={(\dxs,0)}]
    \ftint{0.68}{0}{\fcb}{1.14}
    \ftint{0.68}{-0.50}{\fcb}{\frh}
    \fgrid{0.68}{-0.50}{1}{1}{\fcb}{\frh}
    \fcolvals{0.68}{-0.12}{\fcb}{{$s$}}
    \ftint{1.42}{0}{\fcd}{1.14}
    \fgrid{1.42}{0}{3}{1}{\fcd}{\frh}
    \fcolvals{1.42}{1.14}{\fcd}{{$z\textstyle\sum_k A_{1k}$},{$z\textstyle\sum_k A_{2k}$},{$z\textstyle\sum_k A_{3k}$}}
    \foldbody
    \fcolvals{0.68}{2.88}{\fcb}{{$q_1$},{$q_2$},{$q_3$}}
    \foldafull
    \draw[flow] (1.35,0.57) -- (1.10,0.57);
    \node[axname, anchor=south, inner sep=0.5pt] at (1.225,0.66) {$+$};
    \draw[grpout] (0.65,-0.53) rectangle (1.05,1.17);
    \draw[grpout] (1.39,-0.03) rectangle (2.67,1.17);
    \node[fdlab, anchor=north] at (0.51,-0.58) {fold into\\accumulator};
    \node[fdlab, anchor=north] at (2.03,-0.58) {output\\correction};
    \node[caselab, anchor=south] at (0.38,3.88) {after fold};
  \end{scope}
  \draw[foldbig] (1.14,0.57) -- (1.51,0.57);
\end{tikzpicture}}
\caption{Affine folding of $C=A(sq+z)$, where $K$ is the reduction axis.}
\label{fig:fold-mechanism}
\end{figure}

Consider $C=A(sq+z)$, where $A\in\mathbb{R}^{M\times K}$, $q\in\mathbb{R}^{K\times N}$ holds quantized values, and $k$ indexes the reduction.
When each row $k$ of $q$ shares a scale $s_k$ and offset $z_k$ across columns $j$, distributing reconstruction through the product gives
\begin{equation}
  \begin{aligned}
    C_{i,j} &= \sum_k A_{i,k}(s_kq_{k,j}+z_k) \\
      &= \underbrace{\sum_k(s_kA_{i,k})q_{k,j}}_{\text{scaled input}}+
         \underbrace{\sum_k z_kA_{i,k}}_{\text{output correction}}.
  \end{aligned}
  \label{eq:operand-fold}
\end{equation}
We call this \emph{operand-fold} (\Cref{fig:fold-operand}): the scale multiplies the other input $A$, and the offset contributes a separate reduction of $A$ without reconstructing individual entries of $q$.

\tikzset{
  mgrid/.style={draw=black!55, line width=0.5pt},
  mframe/.style={draw=black!70, line width=0.6pt},
  mname/.style={font=\footnotesize\sffamily, text=black!75},
  fcell/.style={font=\scriptsize, text=black!85, inner sep=0pt},
  axname/.style={font=\footnotesize\sffamily, text=black!65},
  axarrow/.style={-{Stealth[length=3.4pt, width=2.8pt]}, draw=black!55, line width=0.5pt},
  szarrow/.style={{Stealth[length=3.2pt, width=2.6pt]}-{Stealth[length=3.2pt, width=2.6pt]}, draw=ptred!85!black, line width=0.7pt},
  szlab/.style={font=\footnotesize\sffamily, text=ptred!70!black, fill=white, inner sep=0.6pt},
  caselab/.style={font=\footnotesize\sffamily, text=black!85, align=center},
  mode/.style={font=\footnotesize\sffamily\bfseries, text=ptblue!80!black, align=center},
  qbox/.style={draw=black!70, line width=0.5pt, rounded corners=2pt, minimum height=0.54cm, minimum width=0.72cm, inner xsep=3pt, align=center, font=\footnotesize\sffamily},
  quant/.style={qbox, double, double distance=1pt, fill=ptcyan!20},
  prod/.style={qbox, fill=white},
  res/.style={qbox, fill=ptgreen!16},
  flow/.style={-{Stealth[length=4.2pt, width=3.2pt]}, draw=black!75, line width=0.7pt},
  foldline/.style={-{Stealth[length=4pt, width=3pt]}, draw=ptblue, line width=0.8pt, dashed},
  grpout/.style={draw=ptred!85!black, line width=1.1pt, rounded corners=1.2pt},
  foldbig/.style={-{Stealth[length=5pt, width=4pt]}, draw=ptblue, line width=1.2pt},
  foldlab/.style={font=\footnotesize\sffamily, text=ptblue!80!black, fill=white, inner sep=1pt, align=center},
}

\providecommand{\foldmgridd}[5]{%
  \foreach \gi in {1,...,#4}{\ifnum\gi<#4
    \draw[mgrid] (#1+\gi*#5,#2) -- (#1+\gi*#5,#2+#3*#5);\fi}
  \foreach \gj in {1,...,#3}{\ifnum\gj<#3
    \draw[mgrid] (#1,#2+\gj*#5) -- (#1+#4*#5,#2+\gj*#5);\fi}
  \draw[mframe] (#1,#2) rectangle ++(#4*#5,#3*#5);}

\newcommand{\foldkvpanel}[2]{%
\begin{tikzpicture}[x=1cm, y=1cm]
  \def\t{0.38}
  \def\my{0.36}
  \useasboundingbox (-0.20,1.80) rectangle (8.04,-0.44);
  \foreach \r [count=\i from 0] in {1,...,3}{\pgfmathsetmacro{\tt}{14+24*\i}
    \fill[ptblue!\tt] (0.46,\my+3*\t-\i*\t-\t) rectangle ++(3*\t,\t);}
  \foldmgridd{0.46}{\my}{3}{3}{\t}
  \node[axname, anchor=east]  at (0.40,\my+1.5*\t) {$L$};
  \node[axname, anchor=south] at (0.46+1.5*\t,\my+3*\t+0.02) {$D$};
  \draw[szarrow] (0.50,\my+1.5*\t) -- (0.46+3*\t-0.04,\my+1.5*\t);
  \node[caselab, anchor=north] at (0.46+1.5*\t,\my-0.08) {token-wise\\$s,z$ span $D$};
  \draw[foldline] (1.72,\my+1.5*\t) -- (1.96,\my+1.5*\t);
  \node[mode, anchor=west] at (2.02,\my+1.5*\t) {#1};
  \foreach \q [count=\j from 0] in {1,...,3}{\pgfmathsetmacro{\tt}{14+24*\j}
    \fill[ptblue!\tt] (4.30+\j*\t,\my) rectangle ++(\t,3*\t);}
  \foldmgridd{4.30}{\my}{3}{3}{\t}
  \node[axname, anchor=south] at (4.30+1.5*\t,\my+3*\t+0.02) {$D$};
  \draw[szarrow] (4.30+1.5*\t,\my+0.04) -- (4.30+1.5*\t,\my+3*\t-0.04);
  \node[caselab, anchor=north] at (4.30+1.5*\t,\my-0.08) {channel-wise\\$s,z$ span $L$};
  \draw[foldline] (5.56,\my+1.5*\t) -- (5.80,\my+1.5*\t);
  \node[mode, anchor=west] at (5.86,\my+1.5*\t) {#2};
\end{tikzpicture}}
\begin{figure}[t]
\centering
\subfloat[$QK^{\!\top}$ reduces over $D$, so token-wise K uses acc-fold and channel-wise K uses operand-fold.\label{fig:fold-kv-qk}]{\foldkvpanel{acc-fold}{operand-fold}}\\[-2pt]
\subfloat[$PV$ reduces over $L$, so token-wise V uses operand-fold and channel-wise V uses acc-fold.\label{fig:fold-kv-pv}]{\foldkvpanel{operand-fold}{acc-fold}}
\caption{Affine folding modes for token-wise and channel-wise K and V tiles.}
\label{fig:fold-kv}
\end{figure}


When each column $j$ of $q$ shares a scale $s_j$ and offset $z_j$ across the reduction index $k$, both terms factor out of the sum:
\begin{equation}
  \begin{aligned}
    C_{i,j} &= \sum_k A_{i,k}(s_jq_{k,j}+z_j) \\
      &= \underbrace{s_j\sum_k A_{i,k}q_{k,j}}_{\text{scaled partial result}}+
         \underbrace{z_j\sum_k A_{i,k}}_{\text{output correction}}.
  \end{aligned}
  \label{eq:acc-fold}
\end{equation}
We call this \emph{acc-fold} (\Cref{fig:fold-acc}): the scale multiplies the accumulated product, and the offset multiplies the corresponding row sum of $A$.
Both \emph{operand-fold} and \emph{acc-fold} preserve the result of applying $sq+z$ element-wise in exact arithmetic.

In attention, K and V arrive as $[L,D]$ tiles, with $L$ cached tokens and head dimension $D$.
Token-wise groups share a scale and offset along $D$ within each token, whereas channel-wise groups share them along $L$ within each channel.
Because $QK^{\!\top}$ reduces over $D$, token-wise K uses acc-fold and channel-wise K uses operand-fold into $Q$ (\Cref{fig:fold-kv-qk}).
Because $PV$ reduces over $L$, token-wise V uses operand-fold into $P$ and channel-wise V uses acc-fold (\Cref{fig:fold-kv-pv}).
\sys{} allows K and V to select their quantization and folding schemes independently according to the cache's quantization policy.
\Cref{fig:fold-pipeline} illustrates where scales and offsets are applied in the attention algorithm.
\sys{} retains online softmax~\cite{dao2023flashdecoding,dao2022flashattention,dao2024flashattention2} and changes only the QK and PV operand paths and their output corrections.

\tikzset{
  mgrid/.style={draw=black!55, line width=0.5pt},
  mframe/.style={draw=black!70, line width=0.6pt},
  mname/.style={font=\footnotesize\sffamily, text=black!75},
  fcell/.style={font=\scriptsize, text=black!85, inner sep=0pt},
  axname/.style={font=\footnotesize\sffamily, text=black!65},
  axarrow/.style={-{Stealth[length=3.4pt, width=2.8pt]}, draw=black!55, line width=0.5pt},
  szarrow/.style={{Stealth[length=3.2pt, width=2.6pt]}-{Stealth[length=3.2pt, width=2.6pt]}, draw=ptred!85!black, line width=0.7pt},
  szlab/.style={font=\footnotesize\sffamily, text=ptred!70!black, fill=white, inner sep=0.6pt},
  caselab/.style={font=\footnotesize\sffamily, text=black!85, align=center},
  mode/.style={font=\footnotesize\sffamily\bfseries, text=ptblue!80!black, align=center},
  qbox/.style={draw=black!70, line width=0.5pt, rounded corners=2pt, minimum height=0.54cm, minimum width=0.72cm, inner xsep=3pt, align=center, font=\footnotesize\sffamily},
  quant/.style={qbox, double, double distance=1pt, fill=ptcyan!20},
  prod/.style={qbox, fill=white},
  res/.style={qbox, fill=ptgreen!16},
  flow/.style={-{Stealth[length=4.2pt, width=3.2pt]}, draw=black!75, line width=0.7pt},
  foldline/.style={-{Stealth[length=4pt, width=3pt]}, draw=ptblue, line width=0.8pt, dashed},
  grpout/.style={draw=ptred!85!black, line width=1.1pt, rounded corners=1.2pt},
  foldbig/.style={-{Stealth[length=5pt, width=4pt]}, draw=ptblue, line width=1.2pt},
  foldlab/.style={font=\footnotesize\sffamily, text=ptblue!80!black, fill=white, inner sep=1pt, align=center},
}

\providecommand{\foldmgridd}[5]{%
  \foreach \gi in {1,...,#4}{\ifnum\gi<#4
    \draw[mgrid] (#1+\gi*#5,#2) -- (#1+\gi*#5,#2+#3*#5);\fi}
  \foreach \gj in {1,...,#3}{\ifnum\gj<#3
    \draw[mgrid] (#1,#2+\gj*#5) -- (#1+#4*#5,#2+\gj*#5);\fi}
  \draw[mframe] (#1,#2) rectangle ++(#4*#5,#3*#5);}

\begin{figure}[t]
\centering
\begin{tikzpicture}[x=1cm, y=1cm, qbox/.append style={minimum height=0.46cm}]
  \node[prod]  (Q)  at (0.95, 2.70) {$Q$};
  \node[prod]  (qk) at (0.95, 2.04) {QK};
  \node[res]   (S)  at (0.95, 1.38) {$S$};
  \node[prod]  (P)  at (0.95, 0.72) {$P$};
  \node[prod]  (pv) at (0.95, 0.06) {PV};
  \node[res]   (O)  at (0.95,-0.60) {$O$};
  \node[quant] (Kq) at (-4.00, 2.04) {$K^{\!\top}$ quantized};
  \node[quant] (Vq) at (-4.00, 0.06) {$V$ quantized};
  \node[prod]  (Ky) at (-1.45, 2.04) {K operand};
  \node[prod]  (Vy) at (-1.45, 0.06) {V operand};
  \draw[flow] (Q) -- (qk);
  \draw[flow] (qk) -- (S);
  \draw[flow] (S) -- node[axname, right, xshift=0.42cm, inner sep=1pt] {softmax} (P);
  \draw[flow] (P) -- (pv);
  \draw[flow] (pv) -- (O);
  \node[axname, anchor=west] at (1.62,1.38) {K corr.};
  \draw[flow] (1.55,1.38) -- (S.east);
  \node[axname, anchor=west] at (1.62,-0.60) {V corr.};
  \draw[flow] (1.55,-0.60) -- (O.east);
  \draw[flow] (Kq) -- (Ky);
  \draw[flow] (Ky) -- (qk);
  \draw[flow] (Vq) -- (Vy);
  \draw[flow] (Vy) -- (pv);
  \foreach \src/\ysrc/\ydst/\lab in {Kq/2.31/2.70/{channel-wise K\\operand-fold}, Kq/1.77/1.38/{token-wise K\\acc-fold}, Vq/0.33/0.72/{token-wise V\\operand-fold}, Vq/-0.21/-0.60/{channel-wise V\\acc-fold}}{
    \draw[foldline] (-4.00,\ysrc) -- (-4.00,\ydst) -- node[foldlab, pos=0.52] {\lab} (0.55,\ydst);}
\end{tikzpicture}
\caption{Affine folding in attention.
Dashed arrows show where each K and V quantization scheme applies its scales.}
\label{fig:fold-pipeline}
\end{figure}


With $h_q$ query heads per KV head and group size $g$, affine folding reduces scalar scale-and-offset operations from one multiply-add per stored value to $2h_q/g$ (\Cref{sec:appendix-head-scaling}).

\subsection{Affine folding spans multiple quantization groups}
\label{sec:design-groups}

Under acc-fold, an unscaled partial product cannot cross a quantization-group boundary.
When a reduction spans multiple groups, the kernel applies each group's scale and offset to its partial product before adding the result to the running accumulator (\Cref{eq:acc-fold}).
We call this partial product \emph{group-local} because it sums only over the values that share one scale and offset.
Operand-fold instead scales the other matrix input before multiplication, so the accumulator can span multiple groups without rescaling; its offset correction still follows \Cref{eq:operand-fold}.

Both the Hopper and Blackwell kernels enforce the same ordering constraints on these corrections, although they assign the work to different warpgroups (\Cref{sec:design-pipelines}).
For operand-fold, the kernel can accumulate offset contributions for the same output separately and add their sum later.
QK corrections must precede softmax, while PV corrections must precede division by the accumulated softmax denominator.

\subsection{Constructing Tensor Core operands from packed integers}
\label{sec:design-unpack}

\sys{} constructs floating-point Tensor Core operands directly from packed integer bit patterns.%
\footnote{The Hopper kernel uses BF16 operands with \texttt{wgmma}.
The Blackwell kernel uses \texttt{tcgen05} with E4M3 inputs for $Q$ and $P$, E4M3 KV operands for INT4, and E2M1 KV operands for INT2 (using the \texttt{kind::f8f6f4} instruction mode~\cite{nvidia2026ptx}).}
Hopper combines each integer value with a fixed BF16 bit pattern, a technique known as mantissa injection~\cite{kim2022elephants,frantar2025marlin,du2026bitdecoding}.
Combining $q\in\{0,\ldots,15\}$ with \texttt{0x4300} via bitwise OR produces the exact BF16 value $y=q+128$.

On Blackwell, \sys{} can convert integer values into floating point more efficiently by \emph{directly placing each integer's bit pattern in the low bits of a floating-point operand and filling the remaining bits with zeros}.
For example, the bit patterns of the INT2 values $0,1,2,3$ represent $0,0.5,1,1.5$ in FP4 (E2M1), giving $y=q/2$.
INT4 similarly yields FP8 (E4M3) operands with $y=q\cdot2^{-9}$~\cite{lin2025marlinw4a8}.

\begin{table}[t]
  \centering
  \scriptsize
  \begin{tabular*}{\columnwidth}{@{\extracolsep{\fill}}lcccc@{}}
    \toprule
    Mapping & $q\to y$ & $y$ & $(\alpha,\beta)$ & ALU/elem. \\
    \midrule
    \multicolumn{5}{@{}l}{\emph{Hopper}} \\
    \cmidrule(lr){1-5}
    INT4$\to$BF16 & $\texttt{0x4300}\,|\,q$ & $q+128$ & $(1,128)$ & 1.00 \\
    INT2$\to$BF16 & $\texttt{0x4300}\,|\,q$ & $q+128$ & $(1,128)$ & 1.00 \\
    \multicolumn{5}{@{}l}{\emph{Blackwell}} \\
    \cmidrule(lr){1-5}
    INT4$\to$E4M3 & $q$ & $q\cdot2^{-9}$ & $(2^{-9},0)$ & 0.38 \\
    INT2$\to$E2M1 & $q$ & $q/2$ & $(1/2,0)$ & 0.19 \\
    \bottomrule
  \end{tabular*}
  \caption{Tensor Core operand mappings and their static ALU costs per stored value.
  The $q{\to}y$ column gives the bit-level construction of $y$ from $q$.}
  \label{tab:operand-construction}
\end{table}

\begin{figure*}[t]
\centering
\subfloat[Packed register layout and bit-level operand construction.\label{fig:dataflow-packing}]{%
\begin{minipage}[t]{0.49\textwidth}
\centering
\includegraphics[width=\linewidth]{assets/fig_dataflow_packing.pdf}
\end{minipage}}%
\hfill
\subfloat[K operand flow from global memory to the Tensor Core.
Tile shapes are logical (tokens $\times$ channels).\label{fig:dataflow-threads}]{%
\begin{minipage}[t]{0.49\textwidth}
\centering
\includegraphics[width=\linewidth]{assets/fig_dataflow_threads.pdf}
\end{minipage}}
\caption{Register packing and global-to-Tensor-Core dataflow for INT2 K in \sys{}.}
\label{fig:dataflow}
\end{figure*}


These constructions produce operands of the form $y=\alpha q+\beta$ ($\alpha\ne0$) whose numerical values differ from the stored integers.
To account for $(\alpha,\beta)$ without extra per-element instructions, we fold both constants into each group's shared scale and offset.
Substituting $q=(y-\beta)/\alpha$ into $x=sq+z$ (\Cref{eq:dequant}) gives
\begin{equation}
  x=s'y+z',\qquad s'=s/\alpha,\qquad z'=z-\beta s'.
  \label{eq:operand-transform}
\end{equation}
Affine folding then applies $s'$ and $z'$ in place of $s$ and $z$, leaving the attention result unchanged in exact arithmetic (\Cref{sec:design-fold}).
$s'$ and $z'$ are precomputed in FP32 when the quantized KV cache is stored in HBM.
\Cref{tab:operand-construction} summarizes the $q\to y$ mappings, the $(\alpha,\beta)$ pairs used for compensation, and the ALU cost of converting each integer KV element into a Tensor Core operand.

\Cref{fig:dataflow-packing} shows how \sys{} packs INT2 K values in each 32-bit register to match these bit-level constructions.
On Hopper, each byte interleaves pairs of 2-bit values so that one bit shift and one masked merge (\texttt{LOP3}) produce two BF16 operands in a 32-bit register.
On Blackwell, each 32-bit register holds 16 two-bit values arranged as eight channel pairs; one bit shift and two masks unpack the register directly into two 32-bit registers holding 16 FP4 operands.
Both kernels use the Tensor Memory Accelerator (TMA) to copy packed tiles from global to shared memory (\Cref{fig:dataflow-threads}).
Hopper consumers build K operands in registers and collectively issue \texttt{wgmma}; Blackwell prepares a shared-memory operand tile that a separate issue thread supplies to Tensor Cores by issuing \texttt{tcgen05}~\cite{nvidia2026ptx}.
\Cref{sec:operand-construction} extends the INT2 K example to the packing and dataflow of INT2 V and INT4 KV operands on both GPUs.

\begin{figure*}[t]
\centering
\subfloat[H100: \texttt{WG0} issues the tile copies using TMA, and two consumer warpgroups alternate even and odd tiles.\label{fig:arch-hopper}]{%
  \includegraphics[width=\textwidth]{assets/fig_arch_timeline_sm90.pdf}}

\subfloat[B200: \texttt{WG0} issues the tile copies and \texttt{tcgen05} matrix instructions, while \texttt{WG1}--\texttt{WG3} prepare KV operands, apply affine folding, and compute softmax on CUDA Cores.\label{fig:arch-blackwell}]{%
  \includegraphics[width=\textwidth]{assets/fig_arch_timeline_sm100.pdf}}
\caption{\sys{}'s Hopper and Blackwell pipelines.
$\mathrm{FP}(X)$ constructs $X$'s floating-point Tensor Core operand from packed integers (\Cref{sec:design-unpack}).
\texttt{corr.}\ denotes the output correction, and \texttt{WG}$x$ \texttt{W[}$y$\texttt{]} denotes warpgroup $x$, warp $y$.
Dashed boxes execute only for the corresponding K or V quantization axis.}
\label{fig:arch-timeline}
\end{figure*}


\subsection{Warp-specialized GPU pipelines}
\label{sec:design-pipelines}

We design separate warp-specialized pipelines for Hopper and Blackwell.
Hopper's \texttt{wgmma} accumulates into registers owned by its issuing warpgroup, whereas Blackwell's \texttt{tcgen05} accumulates into tensor memory (TMEM) accessible to other warpgroups in the CTA~\cite{nvidia2026ptx}.

On Hopper, one producer warpgroup issues TMA copies, and two consumer warpgroups process alternating even and odd KV tiles (\Cref{fig:arch-hopper}).
Each consumer builds KV operands and runs its own QK--softmax--PV chain, keeping online-softmax state and output accumulators in registers.
Each consumer applies scales and offset corrections according to \Cref{fig:fold-pipeline}: K corrections precede softmax, and V corrections update the output as described in \Cref{sec:design-groups}.
For channel-wise K, the producer warpgroup prepares the scaled $Q$ inputs.
At the end of the sequence, one consumer merges the two warpgroups' softmax states and partial outputs into the final CTA output.

Blackwell assigns copy and matrix-product issue, operand preparation, softmax, and output accumulation to separate warpgroups (\Cref{fig:arch-blackwell}).
The issue warpgroup starts TMA transfers and \texttt{tcgen05} products, while the preparation warpgroup constructs KV operands in shared memory.
The softmax warpgroup reads QK accumulators from TMEM, applies K offset corrections (and K scales for token-wise K) in FP32, computes the softmax weights $P$, and scales $P$ for token-wise V (quantizing folded $Q$ and $P$ tiles to E4M3 with a per-row scale).
The output warpgroup reads PV accumulators from TMEM, maintains the running output, and applies V offset corrections (and V scales for channel-wise V); for channel-wise K, it also prepares the scaled $Q$ inputs.
Shared-memory buffers pass operands between warpgroups, while TMEM passes matrix-product accumulators.

Multi-stage buffers in shared memory and TMEM allow producer warpgroups to work on later tiles while downstream warpgroups process earlier tiles.
We tune the buffer depth and producer lead distance empirically for each workload; for clarity, \Cref{fig:arch-blackwell} illustrates a one-step lead.
